\documentclass[supercite]{Experimental_Report}
\title{深~度~中~国}
\author{阮云轩}
%\coauthor{张三、李四}
\school{计算机科学与技术学院}
\classnum{CS2404}
\stunum{U202490035}
%\costunum{U202115631、U202115631}
\instructor{多教师任课} % 该系列实验报告模板由华科大计院教师陈加忠制作
\date{2025年04月14日}

\usepackage{algorithm, multirow}
\usepackage{algpseudocode}
\usepackage{amsmath}
\usepackage{amsthm}
\usepackage{framed}
\usepackage{mathtools}
\usepackage{subcaption}
\usepackage{xltxtra} %提供了针对XeTeX的改进并且加入了XeTeX的LOGO, 自动调用xunicode宏包(提供Unicode字符宏)
\usepackage{bm}
\usepackage{tikz}
\usepackage{tikzscale}
\usepackage{pgfplots}
%\usepackage{enumerate}
\usepackage{hyperref}

\pgfplotsset{compat=1.16}

\newcommand{\cfig}[3]{
	\begin{figure}[htb]
		\centering
		\includegraphics[width=#2\textwidth]{images/#1.tikz}
		\caption{#3}
		\label{fig:#1}
	\end{figure}
}

\newcommand{\sfig}[3]{
	\begin{subfigure}[b]{#2\textwidth}
		\includegraphics[width=\textwidth]{images/#1.tikz}
		\caption{#3}
		\label{fig:#1}
	\end{subfigure}
}

\newcommand{\xfig}[3]{
	\begin{figure}[htb]
		\centering
		#3
		\caption{#2}
		\label{fig:#1}
	\end{figure}
}

\newcommand{\rfig}[1]{\autoref{fig:#1}}
\newcommand{\ralg}[1]{\autoref{alg:#1}}
\newcommand{\rthm}[1]{\autoref{thm:#1}}
\newcommand{\rlem}[1]{\autoref{lem:#1}}
\newcommand{\reqn}[1]{\autoref{eqn:#1}}
\newcommand{\rtbl}[1]{\autoref{tbl:#1}}

\algnewcommand\Null{\textsc{null }}
\algnewcommand\algorithmicinput{\textbf{Input:}}
\algnewcommand\Input{\item[\algorithmicinput]}
\algnewcommand\algorithmicoutput{\textbf{Output:}}
\algnewcommand\Output{\item[\algorithmicoutput]}
\algnewcommand\algorithmicbreak{\textbf{break}}
\algnewcommand\Break{\algorithmicbreak}
\algnewcommand\algorithmiccontinue{\textbf{continue}}
\algnewcommand\Continue{\algorithmiccontinue}
\algnewcommand{\LeftCom}[1]{\State $\triangleright$ #1}

\newtheorem{thm}{定理}[section]
\newtheorem{lem}{引理}[section]

\colorlet{shadecolor}{black!15}

\theoremstyle{definition}
\newtheorem{alg}{算法}[section]

\def\thmautorefname~#1\null{定理~#1~\null}
\def\lemautorefname~#1\null{引理~#1~\null}
\def\algautorefname~#1\null{算法~#1~\null}

\begin{document}
	
	\maketitle
	
	\clearpage
	
	\pagenumbering{Roman}
	
	\tableofcontents[level=2]
	
	\clearpage
	
	\pagenumbering{arabic}
	
	\section{前言}
	
	一个国家的经营是需要考虑许多因素的，这些因素之多之细是我无法估量的，因此在国家政策上表面上的不足或缺失或许己经是对问题当前最好的处理办法。有时或许人们认为针对这个问题明明很明显有更好的解决办法。但我也说不太清，但我认为起码我能想到的方法别人应该或理应也能想到，所以本文提出或注意到的一些我认为之改善之处或许不过是我的一昧浅见。\\另外，有或多的问题根本是无法解决的，这些问题一般是世界发展引申存在或社会道德无法绝对公立和国家传统文化习俗演变而构成的，例如何谓文化传统，何谓封建陋习这个东西在我的价值观来说根本无法绝对，你认为自分主观区分就行。\\因此，要说一个国家被完美地执行是十分艰巨的，所以评价一个国家运行得是否合理，是好是坏只能相对的述概。面对一些国家现象以我能力我或许只能阐述出来而给不了相关建议。
	\\再次强调本论文无可避免地具有主观性，还请读者注意甄别。
	\\不过在文中我相信你也可以感受到我的价值观，也观迎探讨批评。
	
	
	\section{国情总体分析}
	这个不是本论题的要点因此我会简略地进行叙述。\\
	中国的核心竞争力方面可以说是世界数一数二，有着和美国对抗拉扯的实力，相对于美国来说，中国本土有着稳定的战略，食物，生产力资源，能够面对真正战争爆发被切割中止资源贸易的环境，但人口过于太多，还需向他国进口后备粮食以万一。而在软实力方面，如奢侈品，轻生产业，娱乐设施等控制国民满意度的手段足以支撑人们幸褔度，也具有丰富的文化底蕴和世界历史奇观。但由于人口过多和教育制度的问题，两者使得中国的教育制度处于极度分化残酷的情况，而我认为千古文化传统及现今社会规则使得中国严重缺失创新科研的实力，在国际学术上的话语权不足。
	
	\section{存在问题/不足之处/现况}
	下面会着重给出一些我认为现今国家或许存在的问题，社会反馈及现况乃给出自己的分析理解，这些问题主要是围绕软执行方面的，而这些现象或问题基本上每个地方都会有，但我认为这些现象相比起其他地方较为严重
	\subsection{男女权益}
	我们都知道，最初的女权主义的要求是和男性一样受到同等的权益对待，但为何近件发展使得这个精神主义变得歧化。
	这个歧化问题现象可谓存在了一段时间了亦较为敏感，在亚洲地区现象频繁，在这方面的教育普及我认为我国还有学习进步的地方，起码应从下一代普及改善
	\subsubsection{问题现况}
	这个问题所引申出的问题就是现今社会男女对立，这个对立的情况对于一个国家的发展是十分影响的。由基本人际关系乃至企业的团体工作等都会造成影响，若放由其发展情况只会越来越差，也正如我所说对于这个问题我认为根本方法是更新新一代所需学习的基本社会认知架构，只有这样才能开始杜绝这一情况的恶化。\\
	以下是我对于这个问题的个人理解，还望海涵，欢迎讨论。
	
	\subsubsection{产生原因}
	产生原因我认为多少和本身悠久文化的「谦让避忌」、「女士优先」有关，从小就被教育「大丈夫要让女孩子」的思想。这个因素使得对政策动刀的话会引起部分人群的不满，即使政策的变动或许是合理的。*我认为谦让不过是美德，分清美德和责任我认为应是一个成年人应具备的能力。\\加上近件来社会互联网的高速发展和大力普及，许多不幸女性遭男性不公平对待的事件被社群大力推广使加害者的身份和男性标签化，不少民众强调女性安全力量及政策推广，这一本意是基于以男女性身体力量差异以望加强处分的行为以作震慑保护的作用，但既然目的明确，那么若法律法规认可了身体的差异性，那么作用的范围应该明确针对身体性的行为，而确立后也无法避免个体特殊情况的出现。但现今社会政策似乎在非物理手段上也对性别处理分化了(\ref{fig2-1})(\ref{fig2-2})(\ref{fig2-3})，这违背了开头的初衷，加剧了失衡这一现象的发生。\\
	\begin{figure}[H]
		\begin{center}
			\includegraphics[scale=0.21]{images/2-1.jpg}
			\caption{相关图片}
			\label{fig2-1}
		\end{center}
	\end{figure}
	\begin{figure}[H]
		\begin{center}
			\includegraphics[scale=0.21]{images/2-2.jpg}
			\caption{相关图片}
			\label{fig2-2}
		\end{center}
	\end{figure}
	\begin{figure}[H]
		\begin{center}
			\includegraphics[scale=0.21]{images/2-3.jpg}
			\caption{相关图片}
			\label{fig2-3}
		\end{center}
	\end{figure}
	这里给出部分网址作考实参考吧，有些现在己经找不到了\\
	\\还有几个比较有争议性的事件可以留意一下，这些事件的性质十分简单处理结果却令人不解，而至于是否是由于性别差异所造成的结果还请自行判断吧，毕竟也没人会故意再构建一个一模一样唯性别不同的事件去测试，也无法避免这几年法规的变动是否会有影响\\
	\href{https://www.163.com/dy/article/HJ0SPMH00553XHB9.html}{网络名梗：小仙女刀人爆1500金币之科普}\\
	摘要：这个事件主要是在车辆行至高速时，同行的一名女生因疲劳而睡着，但是精神紧绷的陈某却误以为是网约车司机给她们下了迷药，于是拿出了自己随身携带的水果刀威胁司机立刻停车并要报警。在这个过程中，不小心用刀划伤司机脖子导致司机于某受伤。\\司机于某于是选择报警，在警察到来了解案情后先由两名警察将两名与案情无关的女生送往机场，由刘晖和王健两名警官在现场进行调节。但由于司机于某情绪比较激动于是现场调节失败两人只好带他们回局里继续调节。\\由于错过了当天的飞机，加上期末回家陈某身上以无余钱，在警察局内心崩溃蹲在地上嚎啕大哭。两位警官得知她的情况后对她进行了安抚，同时刘晖中队长慷慨解囊给予她1500元作为机票费用，还答应送她去机场，陈某当场感到的热泪盈眶。\\
	
	\href{https://baike.baidu.com/item/2%C2%B76%E9%95%BF%E6%B2%99%E5%A5%B3%E5%AD%90%E8%B7%B3%E8%BD%A6%E4%BA%8B%E4%BB%B6/56105482}{2·6长沙女子跳车事件}\\
		摘要：2021年2月6日晚9点，湖南长沙年仅23岁的女孩车莎莎从货拉拉车上“跳窗”，后经抢救无效不幸离世 。
		2021年2月21日晚，一条“23岁女生在货拉拉车上跳窗身亡”的新闻在网络传播并引发大量关注。21日深夜，货拉拉方面发布声明称，正配合警方调查。2月22日晚间，家属透露，警方告知现场没有明显紧急刹车痕迹，这成为一个新的疑点。2月23日，长沙警方表示已成立专案组。司机涉嫌过失致人死亡罪被刑拘。3月3日，长沙市高新区公安分局发布关于周某春涉嫌过失致人死亡案件的情况通报。\\
		
		除了传统文化的相关之外，另一助长这一风气的因素是有一部分恶劣的人群打着「为女性发声」之口号，大力要求在不同方面之「特权」，其目的是为了享受社会带来的便利和优待，这类人群往往容易助长恶风，令那些从小成长的孩童们以为这就是社会的标准，然后在长大后継承了这种风气不断恶循环。\\简之，我们追求的是男女平权的社会，我们也不应把事件标签化认为全部坏人都是男性，女性都是被欺负的一方，这个社会不是这样的。事实上，就我自认为中肯地看国家己经给予女性相当的优待了(如女性不构成强奸罪)，这些政策能够实行其实我认为有酌情的考虑，即法律己经对各位让步了那么就应该相互尊重而不是一味地寻求「特权」。\\
		要彻底解决这个问题基本上是不可能的，只能从新一代的基层教育着手，或者执行类似理想国——乌托邦般的国家，所有的法规被完美地执行，不容一丝抗议，但显然这个理想是不能实现的因为国民生活太压抑会带多更多的问题，因此只能寻求相对平衡的社会形态
		
		\subsection{社会资源分配}
		这是我所注意到的另一问题，这个问题往往会发生在一些人口大国中，人口基数大了管理也会变得更加困难，但仍有可参考之地方\\
		为何芬兰，瑞士等国家能有如此高的国民幸褔度，这很大程度和国家的运行操作有关。(\ref{fig2-4})，如何体现一个国家的国民是否幸褔﹖这与国家是否具有顶尖实力无太大关系，一个国家的国民若幸褔的话国民应具有高收入稳定，无环境压力，国民国家心系一起这三大重要要素，如之前发生一件事是关于瑞士“全民发钱”(见参考)之公投提议但被否决了，面对这一简单直接纯利益的提议之所以被拒绝很大程度和本身国家的幸褔度和教育理念有关，这一事件惊动了国际，展示了高度的国民一心整体发展，在这一方面我相信各国都有可学习之地。
		
		\begin{figure}[H]
			\begin{center}
				
\includegraphics[scale=0.30]{images/2-4.jpg}
				\caption{相关图片}
				\label{fig2-4}
			\end{center}
		\end{figure}
		\subsubsection{问题现况及产生原因}
		如今国内社会目对的情景是人口太多所引致的高竞争力和生活高压，同时社会人群提出申请贫困资助的条件过于苛刻，造成的现象是生活水平本身很不是太顺利，但又没有逹到申请条件，处于不上不下的状态，比如说申请低保的条件一般需要当前家庭成员人均收入全部低于当地低保标准和其他后置条件才有资格申请资助。想象一下一个出外打拼的打工人，在外省每月拿4000元的薪水，一个人在外地需要独自面对衣食住行等问题，单是房租水电交通饮食这些必需消费己经基本用剩不了多少了。但一个人的收入却又在社保条件之上，不符合社保的资格。使得他们需要每天高强度地社会活动，失去了作为一个人所必需的基本娱乐消费，沦为了单纯的生产力工具，而像这样的情况在全国有无数的例子。
		\\下面是一些关系社会人们的生活质量，具有时效性供参考。
		\\2023 年《都市上班族睡眠状况报告》显示 “78% 的都市‘打工人’每天睡不够 8 小时，平均睡眠时长为 7.5 小时，睡眠不足 6 小时的人占比 25%，47% 上班族反映睡眠质量不佳，仅有 30.6% 上班族深度睡眠时间达标”。
		\\“农民工群体居住条件相对较差，50.6% 的农民工住在配套不全的单间，16.2% 栖身工棚，人均居住面积不足 5 平方米” 等居住条件相关内容，部分来源于一些社会调查和研究报告，但具体出处未明确提及。
		2022 年进城农民工人均居住面积等数据来源于《2022 年农民工监测调查报告》，报告显示 “2022 年进城农民工人均居住面积 22.6 平方米，比上年增加 0.9 平方米。分城市规模看，不同人口规模的城市农民工人均居住面积均有增加。其中，50 万人以下的城市农民工人均居住面积 27.5 平方米，比上年增加 2.0 平方米；500 万人以上的城市农民工人均居住面积为 17.6 平方米，比上年增加 0.6 平方米”
		\\2024 年全国农民工人均月收入数据来源于 2 月 28 日国家统计局网站发布的 2024 年国民经济和社会发展统计公报，“2024 年末全国农民工总量 29973 万人，比上年增长 0.7%。全国农民工人均月收入 4961 元，比上年增长 3.8%”。
		\\2022 年农民工月均收入及分区域、分行业收入数据来源于《2022 年农民工监测调查报告》，如 “农民工月均收入 4615 元，比上年增加 183 元，增长 4.1%。其中，外出农民工月均收入 5240 元，比上年增加 227 元，增长 4.5%；本地农民工月均收入 4026 元，比上年增加 148 元，增长 3.8%”。
		\\2022-2023 年度《中国美好生活大调查》发现，中国人工作日平均休闲时间为 2.76 小时，比 2021 年少了 7.2 分钟。从不同类型城市看，工作日每天休闲时间不足一小时的人群，二线城市占比最高，达到了 14.38%。
		\subsubsection{改善方法}
		对于这个问题基实把低保条件调低的话(即把收入下限提高)其实不是根本之法，因为总会有刚刚好不合申请条件的人出现，不讨论其他情况或因素，我认为仅针对这个问题较好的做法是把收入水平提高下限再分级处理，并设立缓冲级，即比如第一阶是平均收入在0-3500中，并设对应低保，第二级是3000-5000，并对应设低保，而3000-3500即是我刚所提对的缓冲级，就个别情况调配，另外，应也分省分当地消费水平考虑，把人均水电饮食交通也计算出来算入考虑。当然其实这个问题涉及无数的方面，并没有所说的简单降低申请格就行，但就目前而言，除了大力扶持乡村建设和留守孤儿外，在城市中命生活高压的打工人我认为我们亦注大力关心支持。
		
		\subsection{教育制度及方针现况及产生原因}
		对于教育制度方面也十样复杂，但就现象而言，亚洲尤其韩日中新加坡等地的学生之高压高强度学习环境尤为明显，与欧美等其他地方的教育环境是不同的，欧美地区讲求愉快中学习，自问我们一直提倡在愉快中学习真的是所谓的「愉快」吗，造成这两极分化的根本原因是思想乃教育理念的极其不同，这也是为何外国出现了无数的出色科学家及伟人，无论是科学或是文学艺术精神方面，在创造性方面我自己认为尚有进步空间。而这个问题的产生根本也很简单，国内人口太多教育资源不足以支撑每位我国青少年都能就读本科课程，人口太多带来的是无数的竞争力，每位学子需要面对的是无数的应届生乃不至应届生斗争，争取那稀少的课程学位，在此等恶劣环境下，残酷的中考高考应运而生，学生以成就学校乃分配更多的资源给学校，学校不断构建压力环境望学生中考取得好成绩，中考高考把无数学子推向人生十字路口\\。而现在问题是这样的教育方法是否我们所需要的，是否为社会所需要的，若果极度高压的环境能够使学生，此处不应说学生，就个人而论而非社会，因为讨论社会整体是无意义的，社会的生活质量上升而你不变，这对你毫无帮助，一个普通人只需要考虑个人就好了。若一个人经历残酷的中考高考学习能够使其收获良多有良好的生活质量，就叫做值得为之努力，但现今社会是否真的如此，无数的毕业生工作不对口，那么又如何解释这个现象?而所谓的「应试」式教学是否真的能提高国民知识水平，我实在不清楚，这样的学习真的可以累积精神财富吗。在我看来，造成这个局面的原因是环境与理念，人口密度高竞争力大的环境加上东方的教育理念，使得这样残酷的考试制度是当下唯有的正解，而为何如今会出现学历不够用的情况，显然不是各大学生精益求精无限求学，是因为地方己经没有足够就业岗位容纳他们，不能就业又不想荒废自己的大学生们自然就唯有継续读下去这一条路，因此为何这几年来国家一直推荐人们读研考研名额增多目的显然是为了吸引学生们读下去而掩盖就业问题，在我看来这个情况只会不断迭代下去，始终有一天会稳不住也无法想象那天会发生甚么事。
		\subsubsection{改善方法}
		这个问题其实会出现在其他亚洲地区中，只是中国的高人口使得这个问题更加明显和强烈，若果不打算改变教育理念的话其实只有这一条路可以走，这也是没有办法的事，那么算是缓计的则是大力与设学校这一方法，然后开新岗位，但开新岗位这个又是一大复杂问题了，尤具在这个时代文科和政府工己容不下这么多人时，根本找不到新的低门坎的工作供给大家就业。总的而说，这的确是中国现在一个十分束手无策的问题，要彻底解决的话就要对民族理念动刀了，是否动刀这些事情也当然不由我决定了。因此到头来只能做个吃瓜群众也同时被随波逐流无力改变。
		
		\section{参考事件}
		这里会给出一些参考的事件或网址，很多材料已经有些年份了因此还请注意甄别\\
		\href{https://www.163.com/dy/article/G9UU6DCF0537A693.html}{1.陌生男子往女孩头上浇奶茶后续！男子已被顶格拘留}\\
		\href{https://www.163.com/dy/article/I7DIVRJB0553QMAX.html}{女大学生“扣菜事件”：干完坏事撒腿就跑！罚款300却很嚣张？}\\
		\href{https://www.thepaper.cn/newsDetail_forward_23175763}{男子目睹前女友与人同居后自杀，法院：前女友补偿男子父母2万}\\
		\href{https://zhuanlan.zhihu.com/p/462641661}{怀孕女子跳河自杀，男友劝说无果自行离去，法院判赔32万！}\\
		\href{https://www.xinhuanet.com/world/2016-06/06/c_129043067.htm}{瑞士公投否决“政府无条件发钱”}\\
		
		\section{总结感悟}
		总的来说这不过是纸上谈兵，只言词组是无法解决环境性质的问题的，这些问题每个国家多少都会面对到，这是没有办法的事情，唯一办法即是找到一种可上位替代的思想主义方法，但我这为尚若真有这种主义，那它一定是会对世界权贵的地位产生动摇的，这对权贵来说当然是万万不可的，因此，本质上是两波人群的对抗，要全世界都开心幸褔就要平等权贵的能力资源，而当触碰权贵们的底线时他们便会对百姓施压。因此本质上这个问题是不能有根治的一天的，除非世上再无权贵。
		\\要我来说，与其把本文章当作讨论解决国情的主意这样太夸张的描述还不如当成笔者对国情的看法和谈吐，当然也不能尽信，要有对于自己的判断和理解。
		\\以上。
		
	\end{document}
